% Part: normal-modal-logic
% Chapter: frame-correspondence
% Section: equivalence-S5

\documentclass[../../../include/open-logic-section]{subfiles}

\begin{document}

\olfileid[id]{nml}{frd}{es5}

\olsection{Relasi Ekuivalensi dan \Log{S5}}

Logika modal \Log{S5} dicirikan sebagai himpunan !!{formula}s yang valid pada
semua kerangka universal, yakni setiap dunia dapat diakses dari setiap dunia,
termasuk dari dirinya sendiri. Dalam keadaan demikian, $\Box$ bersesuaian
dengan keniscayaan dan $\Diamond$ dengan kemungkinan: $\Box !A$ benar jika
$!A$ benar pada \emph{setiap} dunia, dan $\Diamond !A$ benar jika $!A$ benar
pada \emph{suatu} dunia. Ternyata, \Log{S5} juga dapat dicirikan sebagai
!!{formula}s yang valid pada semua kerangka refleksif, simetris, dan transitif,
yakni pada semua \emph{relasi ekuivalensi}.

\begin{defn}
  Suatu relasi biner~$R$ pada~$W$ merupakan \emph{relasi ekuivalensi} jika dan
  hanya jika relasi itu refleksif, simetris, dan transitif. Suatu relasi~$R$
  pada~$W$ \emph{universal} jika dan hanya jika $Ruv$ bagi semua $u,v\in W$.
\end{defn}

Karena \Ax{T}, \Ax{B}, dan~\Ax{4} mencirikan kerangka refleksif, simetris, dan
transitif, kerangka yang relasi aksesibilitasnya merupakan relasi ekuivalensi
tepat merupakan kerangka tempat ketiga !!{formula} itu valid. Ternyata relasi
ekuivalensi juga dapat dicirikan oleh kombinasi !!{formula}s lainnya, sebab
kondisi-kondisi yang kita gunakan untuk mendefinisikan relasi ekuivalensi
ekuivalen dengan kombinasi kondisi lain yang sudah dikenal pada~$R$.

\begin{prop}\ollabel{prop:equivalences}
  Kondisi-kondisi berikut ekuivalen:
  \begin{enumerate}
  \item $R$ merupakan relasi ekuivalensi;
  \item $R$ refleksif dan Euklidean;
  \item $R$ serial, simetris, dan Euklidean;
  \item $R$ serial, simetris, dan transitif.
  \end{enumerate}
\end{prop}

\begin{proof}
  Latihan.
\end{proof}

\begin{prob}
  Buktikan \olref[nml][frd][es5]{prop:equivalences} dengan menunjukkan:
  \begin{enumerate}
  \item Jika $R$ simetris dan transitif, maka $R$ Euklidean.
  \item Jika $R$ refleksif, maka $R$ serial.
  \item Jika $R$ refleksif dan Euklidean, maka $R$ simetris.
  \item Jika $R$ simetris dan Euklidean, maka $R$ transitif.
  \item Jika $R$ serial, simetris, dan transitif, maka $R$ refleksif.
  \end{enumerate}
  Jelaskan mengapa hal ini cukup untuk membuktikan bahwa kondisi-kondisi
  tersebut ekuivalen.
\end{prob}

\olref{prop:equivalences} merupakan padanan semantik bagi
\olref[prf][prs]{prop:S5}: proposisi itu memberikan pencirian ekuivalen bagi
logika modal kerangka tempat~$R$ merupakan relasi ekuivalensi (logika yang
secara tradisional disebut~\Log{S5}).

Apa hubungan antara relasi universal dan relasi ekuivalensi? Meskipun setiap
relasi universal merupakan relasi ekuivalensi, jelas bahwa tidak setiap relasi
ekuivalensi universal. Akan tetapi, !!{formula}s yang valid pada semua relasi
universal persis sama dengan formula-formula yang valid pada semua relasi
ekuivalensi.

\begin{prop}
  Misalkan $R$ suatu relasi ekuivalensi, dan bagi setiap $w\in W$ definisikan
  \emph{kelas ekuivalensi} dari~$w$ sebagai himpunan
  $[w]=\{w'\in W:Rww'\}$. Maka:
  \begin{enumerate}
  \item $w\in[w]$;
  \item $R$ universal pada setiap kelas ekuivalensi~$[w]$;
  \item Koleksi kelas-kelas ekuivalensi mempartisi~$W$ menjadi himpunan bagian
    yang saling lepas dan bersama-sama mencakup seluruh~$W$.
  \end{enumerate}
\end{prop}

\begin{prop}\ollabel{prop:S5=univ}
  !!^a{formula}~$!A$ valid pada semua kerangka
  $\mModel{F}=\tuple{W,R}$ tempat~$R$ merupakan relasi ekuivalensi, jika dan
  hanya jika formula itu valid pada semua kerangka
  $\mModel{F}=\tuple{W,R}$ tempat~$R$ universal. Karena itu, logika kerangka
  universal tidak lain dari~\Log{S5}.
\end{prop}

\begin{proof}
  Mudah diperiksa bahwa relasi universal~$R$ pada~$W$ merupakan relasi
  ekuivalensi. Karena itu, jika $!A$ valid pada semua kerangka tempat~$R$
  merupakan relasi ekuivalensi, maka $!A$ valid pada semua kerangka universal.
  Bagi arah lainnya, kita berargumen secara kontrapositif: andaikan $!B$
  merupakan !!a{formula} yang gagal pada suatu dunia~$w$ dalam model
  $\mModel{M}=\tuple{W,R,V}$ yang berbasis pada kerangka~$\tuple{W,R}$,
  dengan~$R$ relasi ekuivalensi pada~$W$. Jadi,
  $\mSat/{M}{!B}[w]$. Definisikan model
  $\mModel{M}'=\tuple{W',R',V'}$ sebagai berikut:
  \begin{enumerate}
  \item $W'=[w]$;
  \item $R'$ universal pada~$W'$;
  \item $V'(p)=V(p)\cap W'$.
  \end{enumerate}
  (Jadi, himpunan dunia~$W'$ dalam~$\mModel{M}'$ digambarkan oleh daerah
  berarsir dalam \olref{fig:partition}.) Mudah dilihat bahwa $R$ dan~$R'$
  berimpit pada~$W'$. Selanjutnya dapat ditunjukkan melalui induksi pada
  !!{formula}s bahwa bagi semua $w'\in W'$ berlaku
  $\mSat{M'}{!A}[w']$ jika dan hanya jika $\mSat{M}{!A}[w']$, bagi setiap
  $!A$ (hal ini masuk akal karena $W'\subseteq W$). Khususnya,
  $\mSat/{M'}{!B}[w]$, sehingga $!B$ gagal dalam suatu model yang berbasis
  pada kerangka universal.
\end{proof}

\begin{figure}[t]
  \centering
  \begin{tikzpicture}[node distance=2cm, auto, thick]
    \clip [rounded corners] (0,0) -- (6,0) -- (6,4) -- (0,4) --  cycle;
    \begin{scope}
      \clip (0,0) -- (2,0)  .. controls (1.5,1.5) and (4.5,1.5)
      .. (4,4) -- (0,4) -- (0,0) -- cycle;
      \filldraw[fill=gray!40] (0,4) -- (0,2) .. controls (1.5,1.5)
      and (4.5,1.5) .. (4.5,0) -- (6,0) -- (6,4) -- (0,4) --  cycle;
    \end{scope}
    \draw [name path=line1] (2,0) .. controls (1.5,1.5) and (4.5,1.5) .. (4,4);
    \draw [name path=line2] (0,2)  .. controls (1.5,1.5) and (4.5,1.5) .. (4.5,0);
    \path [name intersections={of=line1 and line2}];
    \draw [rounded corners, use as bounding box, very thick] (0,0)
    -- (6,0) -- (6,4) -- (0,4) --  cycle;
    \node at (2,3) {$[w]$} ;
    \node at (1,1) {$[u]$} ;
    \node at (3,0.75) {$[v]$} ;
    \node at (4.75,2) {$[z]$} ;
  \end{tikzpicture}
  \caption{Partisi~$W$ menjadi kelas-kelas ekuivalensi.}
  \ollabel{fig:partition}
\end{figure}

\end{document}

